\documentclass[12 pt, letterpaper]{hmcpset}
\usepackage{hw-preamble}

\name{Jayden Thadani}
\class{MATH 142 --- Differential Geometry}
\assignment{Homework assignment \#7}
\duedate{22nd March 2024}

\begin{document}

\begin{problem}[Problem 1.]
	Assume that $F = C \circ T_{\mathbf{a}}$ is an isometry which is the composition of an orthogonal transformation $C$ and a translation $T_{\mathbf{a}} : \mathbf{q} \mapsto \mathbf{q} + \mathbf{a}$. Show that $F = T_{\mathbf{b}} \circ C$ where $T_{\mathbf{b}} : \mathbf{p} \mapsto \mathbf{p} + \mathbf{b}$ is translation by $\mathbf{b} = C(\mathbf{a})$
\end{problem}

\begin{solution}
	We know by the classification of isometries that we can uniquely decompose $F = T \circ D$  where $T$  is a translation, and $D \in O_{3}$ . Specifically, we know that $T$  is given by $T : \mathbf{p} \mapsto \mathbf{p} + F(0)$ and $D = F - F(0)$. We can compute these as
	\[
		\begin{split}
			F(0) & = (C \circ T_{\mathbf{a}})(0) \\
			     & = C(T_{\mathbf{a}} + 0) \\
			     & = C(\mathbf{a})
		\end{split}
	\]
	Thus, $T$  is given by $T : \mathbf{p} \mapsto \mathbf{p} + C(\mathbf{a})$  as desired. That is, $T = T_{\mathbf{b}}$. Then
	\[
		\begin{split}
			D(\mathbf{p}) & = F(\mathbf{p}) - F(0) \\
				      & = C(T_{\mathbf{a}}(\mathbf{p})) - C(\mathbf{a}) \\
				      & = C(\mathbf{p} + \mathbf{a}) - C(\mathbf{a})
		\end{split}
	\]
	which we can see by the linearity of $C$ just gives us
	\[
		D(\mathbf{p}) = C(\mathbf{p}).
	\]
	Thus, we have $D = C$ and $F = T_{\mathbf{b}} \circ C$ as desired.
\end{solution}

\pagebreak

\begin{problem}[Problem 2.]
	Consider the quantities
	\[
		C : \begin{bmatrix}
			x_{1} \\
			x_{2} \\
			x_{3}
		\end{bmatrix} \mapsto
		\frac{1}{3}
		\begin{bmatrix}
			-2 & 2 & -1 \\
			2 & 1 & -2 \\
			1 & 2 & 2
		\end{bmatrix}
		\begin{bmatrix}
			x_{1} \\
			x_{2} \\
			x_{3}
		\end{bmatrix}, \quad \quad
		\mathbf{p} = \begin{bmatrix}
			3 \\
			1 \\
			-6
		\end{bmatrix}, \quad \text{and} \quad
		\mathbf{q} = \begin{bmatrix}
			1 \\
			0 \\
			3
		\end{bmatrix}
	\]
	\begin{enumerate}
		\item Compute $C(\mathbf{p})$ and $C(\mathbf{q})$,
		\item Check that $C(\mathbf{p}) \cdot C(\mathbf{q}) = \mathbf{p} \cdot \mathbf{q}$.
		\item Show that $C$ is \textit{orthogonal} --- that is, $\lVert C(\mathbf{x}) - C(\mathbf{y}) \rVert = \lVert \mathbf{x} - \mathbf{y} \rVert$ and $C(0) = 0$.
	\end{enumerate}
\end{problem}

\begin{solution}
	\begin{enumerate}
		\item We can compute these directly.
			\[
				\begin{split}
					C(\mathbf{p}) & = \frac{1}{3} \begin{bmatrix}
			-2 & 2 & -1 \\
			2 & 1 & -2 \\
			1 & 2 & 2
		\end{bmatrix} \begin{bmatrix}
			3 \\
			1 \\
			-6
		\end{bmatrix} \\
						      & = \frac{1}{3} \begin{bmatrix}
						      	2 \\
							19 \\
							-7
						      \end{bmatrix}
				\end{split}
			\]
			and
			\[
				\begin{split}
					C(\mathbf{q}) & = \frac{1}{3} \begin{bmatrix}
			-2 & 2 & -1 \\
			2 & 1 & -2 \\
			1 & 2 & 2
		\end{bmatrix} \begin{bmatrix}
			1 \\
			0 \\
			3
		\end{bmatrix} \\
						      & = \frac{1}{3} \begin{bmatrix}
						      	-5 \\
							-4 \\
							7
						      \end{bmatrix}
				\end{split}
			\]
		\item Again, we can just compute directly ---
			\[
				\begin{split}
					\mathbf{p} \cdot \mathbf{q} & = \begin{bmatrix}
						3 \\
						1 \\
						-6
					\end{bmatrix} \cdot \begin{bmatrix}
						1 \\
						0 \\
						3
					\end{bmatrix} \\
								    & = 3 - 18 \\
								    & = -15.
				\end{split}
			\]
			and
			\[
				\begin{split}
					C(\mathbf{p}) \cdot C(\mathbf{q}) & = \frac{1}{9} \begin{bmatrix}
						2 \\
						19 \\
						-7
					\end{bmatrix} \cdot \begin{bmatrix}
						-5 \\
						-4 \\
						7
					\end{bmatrix} \\
									  & = \frac{-10 - 76 - 49}{9} \\
									  & = -\frac{135}{9} \\
									  & = - 15.
				\end{split}
			\]
			Thus, we have $C(\mathbf{p}) \cdot C(\mathbf{q}) = \mathbf{p} \cdot \mathbf{q}$.
		\item Since $C$ is linear, we have $C(0) = 0$. Using the property we proved in part (b), we also have that
			\[
				\begin{split}
					\lVert C(\mathbf{x}) - C(\mathbf{y}) \rVert^{2} & = \lVert C(\mathbf{x} - \mathbf{y}) \rVert^{2} \\
											& = C(\mathbf{x} - \mathbf{y}) \cdot C(\mathbf{x} - \mathbf{y}) \\
											& = (\mathbf{x} - \mathbf{y}) \cdot (\mathbf{x} - \mathbf{y}) \\
											& = \lVert \mathbf{x} - \mathbf{y} \rVert^{2}.
				\end{split}
			\] \\
			Then, since both norms have the same sign, we must have
			\[
				\lVert C(\mathbf{x}) - C(\mathbf{y}) \rVert = \lVert \mathbf{x} - \mathbf{y} \rVert
			\]
			as desired. 

			Thus, $C$ is orthogonal.
	\end{enumerate}
\end{solution}

\pagebreak

\begin{problem}[Problem 3.]
	Recall that, for any isogeny $F : \mathbb{R}^{3} \to \mathbb{R}^{3}$ , there exists a unique orthogonal transformation $C$ and translation $T$  such that $F = T \circ C$ . In each case below, decide whether $F$  is an isometry of $\mathbb{R}^{3}$. If so, find its translation and orthogonal parts.
	\begin{enumerate}
		\item $F(\mathbf{p}) = - \mathbf{p}$.
		\item $F(\mathbf{p}) = (p_{3} - 1, p_{2} - 2, p_{1} - 3)$.
		\item $F(\mathbf{p}) = (\mathbf{p} \cdot \mathbf{u}) \mathbf{u}$ where $\lVert \mathbf{u} \rVert = 1$.
		\item $F(\mathbf{p}) = (p_{1}, p_{2}, 1)$.
	\end{enumerate}
\end{problem}

\begin{solution}
	\begin{enumerate}
		\item 
			\[
				\boxed{
					\begin{gathered}
						F \in \mathcal{E}_{3} \\
						C = F \\
						T = I
					\end{gathered}
				}
			\] 

			We know that $F$  is an isometry since
			\[
				\begin{split}
					\lVert F(\mathbf{p}) - F(\mathbf{q}) \rVert & = \lVert -\mathbf{p} - (-\mathbf{q}) \rVert \\
										    & = \lVert -(\mathbf{p} - \mathbf{q}) \rVert \\
										    & = \lVert \mathbf{p} - \mathbf{q} \rVert.
				\end{split}
			\]
			We can also see that $F$  is orthogonal since $F(0) = 0$ . Then $T$  must be given by $T : \mathbf{p} \mapsto 0 + \mathbf{p}$ or, more succinctly, $T = I$ , the identity on $\mathbb{R}^{3}$ .
		\item
			\[
				\boxed{
					\begin{gathered}
						F \in \mathcal{E}_{3} \\
						C : (p_{1}, p_{2}, p_{3}) \mapsto (p_{3}, p_{2}, p_{1}) \\
						T : \mathbf{p} \mapsto (-1, -2, -3) + \mathbf{p}.
					\end{gathered}
				}
			\] 

			We can just note that the composition $F = T \circ C$ does hold, $C$ is an orthogonal transformation (proof below) and $T$ is clearly a translation, so $F$ is the composition of isometries, and itself an isometry.
			\[
				\begin{split}
					\lVert C(\mathbf{x}) - C(\mathbf{y}) \rVert & = \lVert (x_{3}, x_{2}, x_{1}) - (y_{3}, y_{2}, y_{1}) \rVert \\
										    & = \sqrt{(x_{3} - y_{3})^{2} + (x_{2} - y_{2})^{2} + (x_{1} - y_{1})^{2}} \\
										    & = \sqrt{(x_{1} - y_{1})^{2} + (x_{2} - y_{2})^{2} + (x_{3} - y_{3})^{2}} \\
										    & = \lVert (x_{1}, x_{2}, x_{3}) - (y_{1}, y_{2}, y_{3}) \rVert \\
										    & = \lVert \mathbf{x} - \mathbf{y} \rVert.
				\end{split}
			\]
			and $C(0, 0, 0)$ is obviously just $(0, 0, 0)$.
		\item
			\[
				\boxed{
					F \notin \mathcal{E}_{3}.
				}
			\] 

			Suppose $\mathbf{v}$ is a nonzero vector orthogonal to $\mathbf{u}$ (such a vector always exists). Then we have
			\[
				\begin{split}
					F(\mathbf{p} + \mathbf{v}) & = ((\mathbf{p} + \mathbf{v}) \cdot \mathbf{u}) \mathbf{u} \\
								   & = (\mathbf{p} \cdot \mathbf{u} + \mathbf{v} \cdot \mathbf{u}) \mathbf{u} \\
								   & = (\mathbf{p} \cdot \mathbf{u}) \mathbf{u}	\\
								   & = F(\mathbf{p}).
				\end{split}
			\] 

			Thus, $\lVert F(\mathbf{p} + \mathbf{v}) - F(\mathbf{p}) \rVert = 0$ even though $\lVert \mathbf{p} + \mathbf{v} - \mathbf{p} \rVert = \lVert \mathbf{v} \rVert$ is nonzero. That is, $F$ is not an isometry.
		\item
			\[
				\boxed{
					F \notin \mathcal{E}_{3}.
				}
			\]
			Notice that $(0, 0, 0)$ and $(0, 0, 1)$ both get sent to $(0, 0, 1)$ under $F$. Thus, $\lVert F(0, 0, 1) - F(0, 0, 0) \rVert = 0$ even though $\lVert (0, 0, 1) - (0, 0, 0) \rVert = 1$. Thus, $F$ is not an isometry.
	\end{enumerate}
\end{solution}

\pagebreak

\begin{problem}[Problem 4.]
	Two tangent vectors $[\mathbf{v}, \mathbf{p}], [\mathbf{w}, \mathbf{q}] \in \mathrm{T}\mathbb{R}^{3}$ are \textit{parallel} if $\mathbf{v} = \mathbf{w}$. If $T$ is a translation, show that for every tangent vector $[\mathbf{v}, \mathbf{p}]$, the tangent vector $T_{*}([\mathbf{v}, \mathbf{p}])$ is parallel to $[\mathbf{v}, \mathbf{p}]$.
\end{problem}

\begin{solution}
	Since all translations are isometries, we can uniquely decompose $T$ into a translation and an orthogonal transformation. The most obvious choice is $T = T \circ I$ where $I$ is the identity on $\mathbb{R}^{3}$. 

	Since the decomposition is uniquem, we know that for any tangent vector $[\mathbf{v}, \mathbf{p}] \in \mathrm{T}_{\mathbf{p}} \mathbb{R}^{3}$ we have
	\[
		\begin{split}
			T_{*}([\mathbf{v}, \mathbf{p}]) & = [I(\mathbf{v}), T(\mathbf{p})] \\
							& = [\mathbf{v}, T(\mathbf{p})].
		\end{split} 
	\]
	and thus, $T_{*}([\mathbf{v}, \mathbf{p}])$ is parallel to $[\mathbf{v}, \mathbf{p}]$.
\end{solution}

\pagebreak

\begin{problem}[Problem 5.]
	Recall that, given a frame $\mathbf{e}_{1}, \mathbf{e}_{2}, \mathbf{e}_{3}$ in $\mathrm{T}_{\mathbf{p}}\mathbb{R}^{3}$ at $\mathbf{p}$, and a frame $\mathbf{f}_{1}, \mathbf{f}_{2}, \mathbf{f}_{3}$ in $\mathrm{T}_{\mathbf{q}}\mathbb{R}^{3}$	at $\mathbf{q}$, there is a unique isogeny $F : \mathbb{R}^{3} \to \mathbb{R}^{3}$ such that $F_{*}(\mathbf{e}_{i}) = \mathbf{f}_{i}$. 

	Given the frame
	\[
		\begin{split}
			\mathbf{e}_{1} & = \left [ \left ( \frac{2}{3}, \frac{2}{3}, \frac{1}{3} \right ), \mathbf{p} \right ] \\
			\mathbf{e}_{2} & = \left [ \left ( -\frac{2}{3}, \frac{1}{3}, \frac{2}{3} \right ), \mathbf{p} \right ] \\
			\mathbf{e}_{3} & = \left [ \left ( \frac{1}{3}, -\frac{2}{3}, \frac{2}{3} \right ), \mathbf{p} \right ]
		\end{split}
	\]
	at $\mathbf{p} = (0, 1, 0)$ and the frame
	\[
		\begin{split}
			\mathbf{f}_{1} & = \left [ \left ( \frac{1}{\sqrt 2}, 0, \frac{1}{\sqrt 2} \right ), \mathbf{q} \right ] \\
			\mathbf{f}_{2} & = \left [ \left ( 0, 1, 0 \right ), \mathbf{q} \right ] \\
			\mathbf{f}_{3} & = \left [ \left ( \frac{1}{\sqrt 2}, 0, -\frac{1}{\sqrt 2} \right ), \mathbf{q} \right ]
		\end{split}
	\]
	at $\mathbf{q} = (3, -1, 0)$ find $\mathbf{a} \in \mathbb{R}^{3}$ and $C \in O_{3}$ such that the isometry $F = T \circ C$ in terms of the translation $T(\mathbf{p}) = \mathbf{p} + \mathbf{a}$ sends the ``$\mathbf{e}$ frame'' to the ``$\mathbf{f}$ frame''.
\end{problem}

\begin{solution}
	\[
		\boxed{
			\begin{gathered}
				C = \frac{1}{3 \sqrt 2} \begin{bmatrix}
				  3 & 0 & 3 \\
				  -2 \sqrt 2 & \sqrt 2 & 2 \sqrt 2 \\
				  1 & 4 & -1
			  \end{bmatrix} \\
				\mathbf{a} = \frac{1}{3} \begin{bmatrix}
				   	9 \\
					- 4 \\
					-2 \sqrt 2
				   \end{bmatrix}
			\end{gathered}
		}
	\]
	Let $\hat{\mathbf{e}}_{i}$ and $\hat{\mathbf{f}}_{i}$ be the vector parts of $\mathbf{e}_{i}$ and $\mathbf{f}_{i}$ respectively, for each $i \in \{ 1, 2, 3 \}$. 

	Let $F = T \circ C$ be the unique decomposition of $F$ into a translation and an orthogonal transformation. Recalling that $F_{*}([\mathbf{v}, \mathbf{p}]) = [C(\mathbf{v}), F(\mathbf{p})]$, this means that we must have $C(\hat{\mathbf{e}}_{i}) = \hat{\mathbf{f}}_{i}$ for all $i \in \{ 1, 2, 3 \}$. 

	Note that since $\hat{\mathbf{e}}_{1}, \hat{\mathbf{e}}_{2}, \hat{\mathbf{e}}_{3}$ and $\hat{\mathbf{f}}_{1}, \hat{\mathbf{f}}_{2}, \hat{\mathbf{f}}_{3}$ are both bases of $\mathbb{R}^{3}$, there is a well defined linear map $C : \hat{\mathbf{e}}_{i} \mapsto \hat{\mathbf{f}}_{i}$. Now we just need to compute the matrix for this in the standard basis over $\mathbb{R}^{3}$ --- $\mathbf{u}_{1}, \mathbf{u}_{2}, \mathbf{u}_{3}$ with values $(1, 0, 0), (0, 1, 0), (0, 0, 1)$.
	We know that the matrix that maps $\mathbf{u}_{i} \mapsto \hat{\mathbf{f}}_{i}$ is just
	\[
		\begin{bmatrix}
			\vdots & \vdots & \vdots \\
			\hat{\mathbf{f}}_{1} & \hat{\mathbf{f}}_{2} & \hat{\mathbf{f}}_{3} \\
			\vdots & \vdots & \vdots
		\end{bmatrix}.
	\]
	To map $\hat{\mathbf{e}}_{i} \mapsto \hat{\mathbf{f}}_{i}$, we just need the map by $\hat{\mathbf{e}}_{i} \mapsto \mathbf{u}_{i}$. This is the inverse of the map by $\mathbf{u}_{i} \mapsto \hat{\mathbf{e}}_{i}$. Since $\hat{\mathbf{e}}_{1}, \hat{\mathbf{e}}_{2}, \hat{\mathbf{e}}_{3}$ is an orthonormal basis, we can just transpose the matrix. That is,
	\[
		\begin{bmatrix}
			\vdots & \vdots & \vdots \\
			\hat{\mathbf{e}}_{1} & \hat{\mathbf{e}}_{2} & \hat{\mathbf{e}}_{3} \\
			\vdots & \vdots & \vdots
		\end{bmatrix}^{-1} = 
		\begin{bmatrix}
			\cdots & \hat{\mathbf{e}}_{1} & \cdots \\
			\cdots & \hat{\mathbf{e}}_{2} & \cdots \\
			\cdots & \hat{\mathbf{e}}_{3} & \cdots
		\end{bmatrix}
	\]
	Thus,
	\[
		\begin{split}
			\begin{bmatrix}
				\vdots & \vdots & \vdots \\
				\hat{\mathbf{f}}_{1} & \hat{\mathbf{f}}_{2} & \hat{\mathbf{f}}_{3} \\
					\vdots & \vdots & \vdots
			\end{bmatrix}
			\begin{bmatrix}
				\cdots & \hat{\mathbf{e}}_{1} & \cdots \\
				\cdots & \hat{\mathbf{e}}_{2} & \cdots \\
				\cdots & \hat{\mathbf{e}}_{3} & \cdots
			\end{bmatrix}
			\begin{bmatrix}
				\vdots \\
				\hat{\mathbf{e}}_{i} \\
				\vdots
			\end{bmatrix} & = 
			\begin{bmatrix}
				\vdots & \vdots & \vdots \\
				\hat{\mathbf{f}}_{1} & \hat{\mathbf{f}}_{2} & \hat{\mathbf{f}}_{3} \\
					\vdots & \vdots & \vdots
			\end{bmatrix}
			\begin{bmatrix}
				\vdots \\
				\mathbf{u}_{i} \\
				\vdots
			\end{bmatrix} \\ 
			& = 
			\begin{bmatrix}
				\vdots \\
				\hat{\mathbf{f}}_{i} \\
				\vdots
			\end{bmatrix}.
		\end{split}
	\]
	Thus, we can compute $C$.
	\[
		\begin{split}
			C & = \begin{bmatrix}
				\vdots & \vdots & \vdots \\
				\hat{\mathbf{f}}_{1} & \hat{\mathbf{f}}_{2} & \hat{\mathbf{f}}_{3} \\
					\vdots & \vdots & \vdots
			\end{bmatrix}
			\begin{bmatrix}
				\cdots & \hat{\mathbf{e}}_{1} & \cdots \\
				\cdots & \hat{\mathbf{e}}_{2} & \cdots \\
				\cdots & \hat{\mathbf{e}}_{3} & \cdots
			\end{bmatrix} \\
			  & = \frac{1}{\sqrt 2} \begin{bmatrix}
				  1 & 0 & 1 \\
				  0 & \sqrt 2 & 0 \\
				  1 & 0 & -1
			  \end{bmatrix} \frac{1}{3}
			  \begin{bmatrix}
				  2 & 2 & 1 \\
				  -2 & 1 & 2 \\
				  1 & -2 & 2
			  \end{bmatrix} \\
			  & = \frac{1}{3 \sqrt 2} \begin{bmatrix}
				  3 & 0 & 3 \\
				  -2 \sqrt 2 & \sqrt 2 & 2 \sqrt 2 \\
				  1 & 4 & -1
			  \end{bmatrix}.
		\end{split}
	\]
	Since we want $F(\mathbf{p}) = \mathbf{q}$, for $T : \mathbf{p} \mapsto \mathbf{p} + \mathbf{a}$. we need to have $\mathbf{q} = T(C(\mathbf{p}))$ and thus, $\mathbf{a} = \mathbf{q} - C(\mathbf{p})$.
	\[
		\begin{split}
			\mathbf{a} & = \begin{bmatrix}
				3 \\
				-1 \\
				0
			\end{bmatrix} - \frac{1}{3 \sqrt 2} \begin{bmatrix}
				  3 & 0 & 3 \\
				  -2 \sqrt 2 & \sqrt 2 & 2 \sqrt 2 \\
				  1 & 4 & -1
			  \end{bmatrix} \begin{bmatrix}
			  	0 \\
				1 \\ 
				0
			  \end{bmatrix} \\
				   & = \begin{bmatrix}
				   	3 \\
					-1 \\
					0
				   \end{bmatrix} - \frac{1}{3 \sqrt 2} \begin{bmatrix}
				   	0 \\
					\sqrt 2 \\
					4
				   \end{bmatrix} \\
				   & = \frac{1}{3} \begin{bmatrix}
				   	9 \\
					- 4 \\
					-2 \sqrt 2
				   \end{bmatrix}.
		\end{split}
	\]
\end{solution}

\end{document}
